\section{Desviaciones respecto a la Parte 1}\label{sec:desviaciones}

Esta sección reúne todo lo que difiere de la Parte 1. Se distinguen tres tipos:
correcciones de errores del texto de la Parte 1, cambios de diseño decididos por
el grupo e hipótesis que la Parte 1 no precisaba. Ninguna cambia la lógica del
modelo salvo las indicadas como cambios de diseño.

\begin{table}[htbp]
\centering
\small
\caption{Correcciones y cambios de diseño respecto a la Parte 1.}
\label{tab:desviaciones}
\begin{tabularx}{\textwidth}{@{}p{3.3cm}p{2.4cm}p{2.4cm}Y@{}}
\toprule
Aspecto & Parte 1 & Parte 2 & Motivo \\
\midrule
\multicolumn{4}{@{}l}{\textit{Correcciones del texto de la Parte 1}}\\
Intensidad de tráfico ρ del perfil anticipado (X2) & 0,65 & 0,81 &
La franja más cargada del perfil anticipado es la de $-45$ a $-30$~min
(24,4~\%, 793~pers/min), no la de $-30$ a $-15$. Con ella, $\rho = 793/978 \approx 0{,}81$.
No cambian el modelo ni los parámetros; los datos lo confirman
(apartado~\ref{sec:rho}). \\
Caso de 2 tornos por puerta & «la cola crece sin parar» & crece durante la punta &
La cola crece sin cota solo mientras dura la punta con $\rho>1$ (de $-30$ a 0~min);
después se vacía y todos entran. \\
\midrule
\multicolumn{4}{@{}l}{\textit{Cambios de diseño decididos por el grupo}}\\
Instante del cierre de puerta (X6) & $t=+20$~min & $t=-30$~min &
P7 pregunta por un cierre «en hora punta», pero en $t=+20$ llega solo el 2~\%
del público. En un estudio previo con 5 réplicas y seis instantes, la diferencia
sin\,$-$\,con redirección en la espera media de los asignados a la puerta cerrada
era de 4~$\pm$~1~s a $+20$ y de 292~$\pm$~24~s a $-30$, la mayor de las
probadas. Se mantienen la puerta (A-1) y la duración (10~min). \\
Acomodadores en X5 & 3 y 6 & 3, 6 y 9 &
Con 3 la cola de conflictos se satura. Se conservan la probabilidad del 1~\% y la
resolución Triangular(1,~3,~10)~min, y se añade un nivel de 9 para estudiar la
acumulación en lugar de ajustar los parámetros. \\
Número de réplicas & unas 30 por escenario & 30 a 247 &
Con el criterio de la propia Parte 1 (IC del 95~\% con semianchura $\le$~5~\% de la
espera media), el piloto estimó tamaños muy distintos. Se usan al menos 30 y el
tamaño necesario de cada escenario, comprobando la precisión al terminar
(apartado~\ref{sec:campana}). \\
\bottomrule
\end{tabularx}
\end{table}

\subsection{Hipótesis adoptadas}

La Parte 1 no precisaba los aspectos de la tabla~\ref{tab:hipotesis}. En cada
uno se ha elegido la lectura más sencilla y fiel al enunciado.

\begin{table}[htbp]
\centering
\small
\caption{Hipótesis adoptadas donde la Parte 1 no era explícita.}
\label{tab:hipotesis}
\begin{tabularx}{\textwidth}{@{}p{4.2cm}Y@{}}
\toprule
Aspecto & Hipótesis adoptada \\
\midrule
Cola formada al cerrar una puerta & Ver el párrafo siguiente. \\
Reloj de las averías & El tiempo entre fallos cuenta solo en funcionamiento;
240~min es el tiempo medio de buen funcionamiento. \\
Recuento de averías & Se cuentan todas las averías registradas en la réplica,
también las posteriores a la entrada. No equivale al número de averías que
perjudicaron a alguien. \\
Relectura y prioridad & La relectura no libera el torno; la prioridad no expulsa un
control en curso. \\
Utilización de los tornos & Se mide con el número nominal de tornos, desde la
apertura hasta el último control. Las averías cuentan como tiempo disponible no
ocupado. \\
Distancia en la evacuación & Quien está caminando hacia su asiento al sonar la
alarma recorre la distancia completa hasta la salida: no se conoce su posición y
es el caso más desfavorable. \\
Cola exterior al sonar la alarma & Quien está fuera, en cola o pasando el torno,
abandona y no cuenta en la evacuación. \\
Salida bloqueada & Se bloquea la salida 1 de la macrozona A; las 6 son iguales,
así que la elección no cambia el resultado. \\
Elección de salida & Se elige al llegar a la zona de salidas, cuando se ven las
colas, y no al levantarse. \\
Conflicto en curso al sonar la alarma & Se da por terminado: el espectador deja el
asiento para evacuar. \\
\bottomrule
\end{tabularx}
\end{table}

\paragraph{Cola formada al cerrar una puerta.}
Hipótesis adoptada: el cierre solo afecta a quien \emph{llega} mientras la puerta
está cerrada; quien ya estaba en su cola o en un torno sigue siendo atendido. Es
la lectura literal de la Parte 1 («nadie entra por esa puerta», «las llegadas se
redirigen o esperan»), que no trata la cola existente. La alternativa, parar
también esa cola, daría un efecto mayor del cierre. Con el cierre en $t=-30$
apenas importa: en el estudio previo había 0 o 2 personas en cola al cerrar,
frente a unas 380 llegadas durante el cierre, así que la hipótesis no influye
en X6. Con cierres más tardíos ($-20$, $-10$) llegaba a haber unas 100 personas
en cola al cerrar; para esos instantes no se ha comprobado y cambiar la
hipótesis obligaría a repetir X6.

Esta hipótesis va unida a una limitación de la redirección: quien llega a la
puerta cerrada conoce al instante la cola de todas las demás, elige la menor y no
camina ninguna distancia extra. La redirección del modelo es, por tanto,
optimista respecto a una real (apartado~\ref{sec:p7}).
